\chapter{Rezultati}\label{cha:rezultati}
V tem poglavju študent predstavi rezultate meritev, analiz, preračunov in drugih obdelav podatkov, pridobljenih v okviru raziskave. Rezultati naj bodo predstavljeni objektivno.

Rezultati morajo biti jasno strukturirani, pregledno prikazani (preglednice in slike) in nedvoumno povezani z metodologijo iz poglavja \ref{cha:metodologija}. V besedilu se poudarijo ključni trendi in značilnosti, ne pa se ponavljajo vse vrednosti iz preglednic in slik. Če je naloga obsežnejša ali razdeljena na več vsebinskih sklopov, se rezultati posameznih sklopov predstavijo v ločenih podpoglavjih, pri čemer mora biti logika razdelitve jasno razvidna.

Poglavji \ref{cha:rezultati} in \ref{cha:diskusija} lahko združite v enotno poglavje \ref{cha:rezultati} Rezultati in diskusija, v kolikor je za izboljšanje jasnosti, razumevanja in berljivosti zaključnega dela to bolj primerno. V tem primeru prikazane ter opisane rezultate sproti interpretirate in o njih razpravljate.
% Po potrebi naslov spremenite v »Rezultati in diskusija« in v glavni
% datoteki odstranite \chapter{Diskusija}\label{cha:diskusija}
V tem delu predstavite vaše razumevanje rezultatov, njihovo interpretacijo in morebitne komentarje. Pri tem se ustrezno sklicujte na podpoglavja, slike in preglednice iz poglavja Rezultati.

Poglavji \ref{cha:rezultati} in \ref{cha:diskusija} lahko združite v enotno poglavje \ref{cha:rezultati} Rezultati in diskusija, v kolikor je za izboljšanje jasnosti, razumevanja in berljivosti zaključnega dela to bolj primerno. V tem primeru prikazane ter opisane rezultate sproti diskutirate.
.
